\scanpage{19}
\mspage{17}

d'où

\[
 f''=f'\circ f:K_\bullet\longrightarrow M_\bullet,\quad\text{et}\quad c'':f''(e_{K_\bullet})\mathbin{\unsure{-|\longrightarrow}}e_{M_\bullet}
\]

défini ainsi

\diagram{\begin{tikzcd}[ampersand replacement=\&,column sep=small]
\deleted{f'(c):} \& f'(f(e_{\unsure{K|k}}))\arrow[r,"{f'(c)}"]\arrow[rr,bend right=14,"{c''\ \text{par déf.}}"'] \& f'(e_L)\arrow[r,"{c'}"] \& e_M .
\end{tikzcd}}

Cela nous donne donc bel et bien accouplements \unsure{dans} $\mathrm{Hot}_{\bullet}$\inserted{(et aussi \deleted{\unsure{dans}} \inserted{dans} $\unsure{\mathrm{Cat\,hot}_{\bullet}|\mathrm{Cat\,hot}_{:}}$ ?)}. Il est difficile de croire qu'ils pourraient ne pas être associatifs !

\margin{?}

D'ailleurs $\mathrm{Hot}_{\bullet}$, avec le produit cartésien ordinaire$^{*}$, satisfait à l'axiomatique \unsure{de|du} \unsure{tantôt|tantôt prévue} pour $\unsure{\mathcal H|\mathrm H}$, et $\mathrm{Hot}_{\bullet}\longrightarrow\mathrm{Hot}$ est compatible avec ce produit$^{**}$, et avec $\unsure{\Pi|\pi}_0$.

% slanted
\margin{$^{*}$ Non, il faut prendre le produit $X_\bullet\mathbin{\unsure{*|\#}}Y_\bullet$, \deleted{\unsure{$(X\times Y)$}} $(X_\bullet\times Y_\bullet,\ X_\bullet\vee Y_\bullet$ contracté en $e$).}

% slanted
\margin{$^{**}$ Non, mais on \unsure{peut} \unsure{faire} jouer quand-même \unsure{un rôle}}

Si tout ceci marche, on aurait donc une théorie des $\mathrm{Hom\,cat}_{\bullet}$ à valeurs dans $\mathrm{Hot}_{\bullet}$ ; et sans doute aussi pour les $\mathrm{CH}_{\infty}$ comme objets de $\deleted{\unsure{\mathrm{Ind}}}\mathrm{Hot}_{\bullet}$, et non seulement de $\mathrm{Ind\,Hot}_{\unsure{\mathrm{eq}|q}}$.

Mais de \deleted{\unsure{toutes}} façons on ne peut en général laisser tomber $\mathrm{Cat\,hot}$, qui reste importante -- donc il faudrait avoir des objets $\mathrm{CH}_{\infty}(a,b)$ de $\mathrm{Ind\,Cat\,hot}_{\bullet}$, pour avoir une \unsure{côte|cote} bien taillée.

Les $\Pi_0(\unsure{\mathrm{Ch}|\mathcal C}(a,b))$ seront à présent des ensembles pointés, et ce sera pour l'image de l'application nulle $a\xrightarrow{0}b$ \unsure{i.e.} le composé

\[
 a\longrightarrow 0\longrightarrow b
\]
% archive sheet number: 18
